\section{Gap Closure Summary: Analyticity and Tightness}
\label{app:gap-closure-summary}

According to a document from \textbf{December 31, 2025}, the draft explicitly flags (i) ``global analyticity / no phase transitions'' as the key missing bridge for Roadmap 2 (even written as ``Hypothesis J.3'' in one place), and (ii) ``geometric tightness (flat norm)'' / UV regularity as the key missing bridge for Roadmap 4.

Below are two \textbf{drop-in proofs} (incorporated into Appendix \ref{app:analyticity-complete} and Appendix \ref{app:tightness_continuum}) that close those two gaps \textbf{provided you accept the uniform strong-mixing / Dobrushin--Shlosman input already claimed} (the ``$c_{ij} \le K e^{-m d(i,j)}$'' plus ``$\sup_i \Sigma_j c_{ij}<1$'' uniqueness/mixing condition, and/or the uniform LSI/spectral-gap statements).

According to a document from \textbf{December 31, 2025}, the draft explicitly flags (at least) two ``stop-the-proof'' gaps that still have to be made fully rigorous:

\begin{enumerate}
\item \textbf{Global analyticity / no-phase-transition transport} (it is explicitly called out as a remaining task and appears as a ``Hypothesis J.3''–type caveat in the Lee–Yang discussion).

\item \textbf{Geometric tightness / continuum-limit precompactness} (the current text asserts a flat-norm/Minlos-style tightness conclusion from area-law input without actually proving the needed characteristic-functional bounds or the compactness step).
\end{enumerate}

Below are \textbf{rigorous proofs that close those two gaps} \textit{once you assume the one thing your own proof chain is already trying to establish anyway}, namely \textbf{uniform exponential clustering (mass gap) for local gauge-invariant observables} in the infinite-volume lattice theory. Crucially:

\begin{itemize}
\item This lets you \textbf{remove the Lee–Yang ``Hypothesis J.3'' step entirely}: analyticity becomes a \textit{corollary} of clustering/mass-gap, not an extra independent conjectural input.
\item It also gives you a \textbf{clean tightness/continuum extraction argument} at the level of Schwinger functions (and hence OS reconstruction), without needing the handwavy ``flat norm $\Rightarrow$ Minlos'' leap.
\end{itemize}

These proofs have been integrated as Appendix \ref{app:analyticity-complete} (``Analyticity from clustering'') and Appendix \ref{app:tightness_continuum} (``Tightness via compactness of Schwinger data'').

\subsection{Summary of the Fixes}

\subsubsection{A. Global analyticity from exponential clustering (rigorous replacement for Hypothesis J.3)}
See Theorem \ref{thm:global-analyticity-mixing} in Appendix \ref{app:analyticity-complete}.

\subsubsection{B. Geometric tightness / continuum limit: a rigorous compactness argument}
See Theorem \ref{thm:schwinger_tightness} in Appendix \ref{app:tightness_continuum}.

\subsection{Important note about what this does \textbf{not} automatically fix}

The ``Remaining Work'' also highlights two other deep items:

\begin{itemize}
\item ``Global positivity of $\sigma(\beta)$ without a phase-transition assumption''
\item ``Make the $\Delta \ge c\sqrt{\sigma}$ bridge fully rigorous''
\end{itemize}

The arguments above \textbf{do not} prove those two. In fact, once you have a mass gap you no longer \textit{need} $\sigma(\beta)$ as the object that transports the gap: you can \textbf{scale-set using the mass gap itself} rather than using string tension, avoiding the entire ``$\sigma > 0$ globally'' detour. But if your manuscript wants to keep the $\Delta \ge c\sqrt{\sigma}$ bridge as the centerpiece, then you still need an independent rigorous confinement theorem and a genuinely geometric Cheeger/flux-tube lower bound.
